\chapter{Resumen del módulo}
\label{sec:resumen}

El módulo de \textit{\Modulo} estudia los fundamentos y la aplicación de mecanismos criptográficos para proteger la información almacenada y en tránsito. Su contenido comprende el cifrado simétrico y asimétrico, las funciones de resumen, la derivación de claves, los esquemas de firma digital y las infraestructuras de clave pública. Estos conceptos se relacionan con los protocolos de comunicación segura, la gestión de identidades y el análisis de debilidades criptográficas.

Por otro lado, la esteganografía complementa este marco mediante el estudio de técnicas que buscan ocultar la presencia de información en un soporte o canal, así como de los procedimientos que permiten detectar esa ocultación.

En conjunto, el módulo proporciona una base para distinguir las propiedades de confidencialidad, integridad y autenticidad, y para valorar su aplicación en sistemas reales.

La práctica trata principalmente los siguientes contenidos: gestión de claves, cifrado autenticado y firma digital mediante herramientas como \textit{OpenSSL}.
